% A \thechapter set in the preamble of a book: lettered chapters, with
% sections, a theorem, and an equation numbered within them, references to
% each with hyperref, the table of contents, and unnumbered front-matter and
% back-matter chapters.
\documentclass{book}
\usepackage{amsmath}
\usepackage{hyperref}
\renewcommand\thechapter{\Alph{chapter}}
\newtheorem{proposition}{Proposition}[chapter]
\begin{document}
\frontmatter
\chapter{Preface}\label{preface}
Front text.
\tableofcontents
\mainmatter
\chapter{Start}\label{start}
\section{Basics}\label{basics}
\begin{proposition}\label{prop}
A claim, see \ref{basics}.
\end{proposition}
\begin{equation}
x = 1 \label{eq}
\end{equation}
Ref: \ref{start}; \ref{basics}; \ref{prop}; \eqref{eq}; \ref{second}.

In italics, \textit{\ref{basics}}; in sans-serif, \textsf{\ref{prop}}; in typewriter, \texttt{\ref{start}}.

In a footnote\footnote{In the note, \ref{basics}.}, too.
\begin{proposition}
In a theorem, \ref{second}.
\end{proposition}
\section{About \ref{start}}
Autoref: \autoref{start}; \autoref{basics}; \autoref{eq}.
\chapter{Second}\label{second}
\section{More}\label{more}
Here \ref{second}; \ref{more}.
\backmatter
\chapter{Afterword}\label{afterword}
Back text.
\end{document}
